% SOURCE_AUTHORITY: sga5_fr_workpass.tex, lines 12396--12402 (LF-normalized unit).
% SOURCE_SHA256: 791F4EFFC5E02832D5D77ED03518C8156D6F07E4C8238B03545DB93D883FBB28
\subsubsection*{3.4}

Sejam $X^\bullet$ e $Y^\bullet$ dos complexos de $A$-módulos. Assim como no caso dos módulos, se um dos complexos $X^\bullet$, $Y^\bullet$ tem uma estrutura adicional de $A$-$B$-bimódulo, é possível introduzir em $\Hom_A^\bullet(X^\bullet,Y^\bullet)$ uma estrutura de complexo de $B$-módulos e, consequentemente, uma estrutura correspondente em $\RHom_A(X^\bullet,Y^\bullet)$.

Suponhamos, por exemplo, que $X^\bullet$ é um objeto de $D^-(A)$ e $Y^\bullet$ é um complexo de $A$-$B$-bimódulos; para precisar, digamos que estamos no caso descrito pelo símbolo ${}_{A-B}Y^\bullet$. Em particular, $Y^\bullet$ é um objeto de $D(A)$. Seja $P^\bullet\to X^\bullet$ uma resolução projetiva de $X^\bullet$. Então $\Hom_A^\bullet(P^\bullet,Y^\bullet)$ é evidentemente um complexo de $B$-módulos pela esquerda, isto é, um objeto de $D(B)$. Se se utiliza outra resolução projetiva $P_1^\bullet\to X^\bullet$, o complexo de $B$-módulos $\Hom_A^\bullet(P_1^\bullet,Y^\bullet)$ é canonicamente isomorfo em $D(B)$ ao complexo $\Hom_A^\bullet(P^\bullet,Y^\bullet)$. Resulta, portanto, que se pode definir $\RHom_A(X^\bullet,Y^\bullet)$ como um objeto de $D(B)$ determinado salvo isomorfismo único, que depende funtorialmente de $X^\bullet$ em $D^-(A)$ e de $Y^\bullet$ em $D(A\otimes_{\mathbb Z}B)$.

Aplicam-se considerações análogas ao complexo $X^\bullet\otimes_A^{\mathbf L}Y^\bullet$. Se, por exemplo, $Y^\bullet$ é um complexo de $A$-$B$-bimódulos como antes e $X^\bullet$ é um complexo de $A$-módulos pela direita limitado à direita, pode-se definir $X^\bullet\otimes_A^{\mathbf L}Y^\bullet$ como um objeto de $D(B)$ determinado salvo isomorfismo único.
